\section{When a joint test can reject what single tests accept}\label{sec:theory}

\begin{lemma}[Composition]\label{lem:composition}
If $P_A$ and $P_B$ are orthogonal projectors with $P_AP_B=0$, then for all inputs $r$, $s$ and $t$,
\[
  J(r,s,t) \;=\; \SA\bigl(\SB(z_r,t),\,s\bigr) \;=\; \SB\bigl(\SA(z_r,s),\,t\bigr).
\]
\end{lemma}

\begin{proof}
$\SA(\SB(z_r,t),s)=\SB(z_r,t)+(z_s-\SB(z_r,t))P_A=z_r+(z_t-z_r)P_B+(z_s-z_r)P_A-(z_t-z_r)P_BP_A$, and the last
term vanishes because the projectors are symmetric, so $P_BP_A=(P_AP_B)^{\top}=0$. Exchanging the roles of $A$ and
$B$ gives the second equality.
\end{proof}

A joint intervention is therefore a single intervention whose recipient is the output of another single
intervention. Whether that recipient is the natural state it simulates decides whether the joint test can tell us
anything new (\cref{fig:identity}). We say that the $B$-edits of an alignment land on their targets if
$\SB(z_r,t)=z_{(a_r,b_t,n_r)}$ for every recipient $r$ and donor $t$.

\begin{proposition}\label{prop:redundant}
If the $B$-edits of an alignment land on their targets, then every joint query $(r,s,t)$ edits the site into the
same state, and makes the same prediction, as the single $A$-query with recipient $r'=(a_r,b_t,n_r)$ and donor
$s$. Hence an alignment that passes all 729 single $A$-queries passes every joint query, and one that passes every
$A$-query of the full-domain single panel passes every query of the joint panel. The same holds with $A$ and $B$
exchanged.
\end{proposition}

\begin{proof}
By \cref{lem:composition}, $J(r,s,t)=\SA(z_{r'},s)$, which is the single $A$-query $(r',s)$. That query predicts
the output of $(a_s,b_{r'},n_{r'})=(a_s,b_t,n_r)$, which is the joint prediction. If the joint query is in the joint
panel, $a_s\ne a_r=a_{r'}$, so the single query is in the full-domain single panel.
\end{proof}

\begin{remark}\label{rem:null}
The proposition survives edits that also move directions the downstream computation ignores. Suppose
$\SB(z_r,t)=z_{r'}+\eta$ with $\eta$ in a subspace $N$ such that $g(w+\nu)=g(w)$ for every $w\in\R^d$ and
$\nu\in N$, and $NP_A\subseteq N$. Then $J(r,s,t)=\SA(z_{r'},s)+\eta(I-P_A)$ with $\eta(I-P_A)\in N$, so the joint
query and the single query $(r',s)$ produce the same output. In Study~1 the coordinates of the nuisance token, which
the network never reads, form such a subspace for the projectors of every valid alignment.
\end{remark}

The proposition assumes that every $B$-edit lands. Applied to one joint query at a time, its argument says when a
joint test can reject what single tests accept.

\begin{corollary}\label{cor:irreducible}
Suppose a joint query $(r,s,t)$ fails while the two single queries it reduces to pass: the $A$-query with
recipient $(a_r,b_t,n_r)$ and donor $s$, and the $B$-query with recipient $(a_s,b_r,n_r)$ and donor $t$. Then both
single edits miss their targets: $\SA(z_r,s)\ne z_{(a_s,b_r,n_r)}$ and $\SB(z_r,t)\ne z_{(a_r,b_t,n_r)}$, also up
to a subspace $N$ as in \cref{rem:null} with $NP_A\subseteq N$ and $NP_B\subseteq N$. In particular, an alignment
that passes all single queries and fails a joint query has, for both variables, single edits that give the
correct output while missing their targets.
\end{corollary}

\begin{proof}
Write $r'=(a_r,b_t,n_r)$ and suppose that $\SB(z_r,t)=z_{r'}+\eta$ with $\eta\in N$ ($\eta=0$ if the $B$-edit lands).
Since $\SA(w,s)=w(I-P_A)+z_sP_A$ for every state $w$, \cref{lem:composition} gives
$J(r,s,t)=\SA(z_{r'}+\eta,s)=\SA(z_{r'},s)+\eta(I-P_A)$, and $\eta(I-P_A)\in N$ because $NP_A\subseteq N$. The joint
query therefore produces the output of the single query $(r',s)$, which predicts the same output, that of
$(a_s,b_t,n_r)$. That single query passes, so the joint query passes too, contrary to the hypothesis. The same
argument with $A$ and $B$
exchanged, using $NP_B\subseteq N$, applies to $\SA(z_r,s)$ and the single query $((a_s,b_r,n_r),t)$. For the last
statement, a failing joint query of an alignment that passes all single queries meets the hypothesis, and its edits
$\SA(z_r,s)$ and $\SB(z_r,t)$ are those of the single queries $(r,s)$ and $(r,t)$, which pass.
\end{proof}

We call a joint failure of this kind \emph{irreducible}: it cannot be traced to a failing single query on a
natural recipient. \citet{grant2026divergent} call intervened states that the network does not reach naturally
\emph{divergent}, and a perturbation that is behaviourally null in some contexts but changes behaviour in others a
\emph{dormant behavioural change}. The single edits of \cref{cor:irreducible} are of this kind: their effect is
null in the single test and appears once the other variable has been edited. Their Theorem~A.2, which shows that a
set of states is closed under patching subsets of coordinates exactly when it is the product of its coordinate
ranges, is the closest formal precedent for the property used in \cref{prop:redundant}. The condition is necessary,
not sufficient: single edits that miss their targets need not combine into a failure.

\begin{figure}[t]
\centering
\begin{tikzpicture}[>=Stealth, font=\small,
    nat/.style={circle, fill=black, inner sep=1.6pt},
    off/.style={circle, draw=black, thick, fill=white, inner sep=1.6pt},
    region/.style={rounded corners=10pt, fill=gray!12, draw=gray!55},
    note/.style={gray!60!black, font=\footnotesize}]
  % (a) B-edits land on natural states
  \begin{scope}
    \draw[region] (0,0) rectangle (5.0,1.5);
    \node[note, anchor=south west] at (0.12,0.08) {natural states};
    \node[nat] (r) at (0.55,0.85) {};
    \node[anchor=east] at (r.west) {$z_r$};
    \node[nat] (rp) at (3.7,0.85) {};
    \node[anchor=north, font=\footnotesize] at ([yshift=-1pt]rp.south) {$z_{(a_r,b_t,n_r)}$};
    \node[off] (j) at (3.7,2.75) {};
    \node[anchor=west] at ([xshift=2pt]j.east) {$J(r,s,t)$};
    \draw[->, thick] (r) -- node[above, font=\footnotesize] {$B$ from $t$} (rp);
    \draw[->, thick] (rp) -- node[right, font=\footnotesize] {$A$ from $s$} (j);
    \node[align=center, font=\footnotesize, anchor=north] at (2.5,-0.25)
      {(a) the $B$-edit lands on its target:\\ the joint query is the single $A$-query $((a_r,b_t,n_r),\,s)$};
  \end{scope}
  % (b) both single edits leave the natural states
  \begin{scope}[xshift=8.2cm]
    \draw[region] (0,0) rectangle (5.0,1.5);
    \node[note, anchor=south east] at (4.9,0.08) {natural states};
    \node[nat] (r) at (1.2,0.75) {};
    \node[anchor=north east] at (r.south west) {$z_r$};
    \node[off] (sa) at (0.45,2.6) {};
    \node[anchor=east] at ([xshift=-2pt]sa.west) {$\SA(z_r,s)$};
    \node[off] (sb) at (4.3,2.15) {};
    \node[anchor=west] at ([xshift=2pt]sb.east) {$\SB(z_r,t)$};
    \node[off] (j) at (2.9,3.25) {};
    \node[anchor=south] at ([yshift=2pt]j.north) {$J(r,s,t)$};
    \draw[->, thick] (r) -- node[left, font=\footnotesize, pos=0.7] {$A$ from $s$} (sa);
    \draw[->, thick] (r) -- node[below right, font=\footnotesize, pos=0.5, inner sep=1pt] {$B$ from $t$} (sb);
    \draw[->, thick, dashed] (sa) -- (j);
    \draw[->, thick, dashed] (sb) -- (j);
    \node[align=center, font=\footnotesize, anchor=north] at (2.5,-0.25)
      {(b) both single edits miss their targets:\\ the joint query is not a single query on a natural recipient};
  \end{scope}
\end{tikzpicture}
\caption{By \cref{lem:composition}, a joint edit is a single edit applied to the result of the other single
edit. (a) If the $B$-edit of $z_r$ lands on its target, the natural state of the input it simulates, the joint
query coincides with an ordinary single query (\cref{prop:redundant}). (b) A joint query can only fail where the
single queries it reduces to pass if both single edits miss their targets (\cref{cor:irreducible}). Filled dots are natural states; open dots are edited states.}
\label{fig:identity}
\end{figure}

\paragraph{The condition can be met.}
A small explicit network, built on the mechanism of the example of \citet[Section~4.2]{grant2026divergent} (a unit
that is silent on every natural state), shows that joint tests can reject an alignment that every single test
accepts. The target is $y=a+b$. The site has four coordinates $z=(u_a,u_b,e_A,e_B)$ with natural value
$(a,b,0,0)$, and the readout is the two-layer ReLU network
\[
  g(z) \;=\; u_a+u_b+\tan\theta\,(e_A+e_B)-k\,\mathrm{relu}\bigl(|e_A|+|e_B|-\tau\bigr),\qquad
  |e|=\mathrm{relu}(e)+\mathrm{relu}(-e),
\]
with $\tau=2.5\sin\theta\cos\theta$ and $k=1/(\sin\theta\cos\theta)$, so that $g(z_x)=a+b$ on every natural state.
The alignment places $A$ on $v_A=(\cos\theta,0,\sin\theta,0)$ and $B$ on $v_B=(0,\cos\theta,0,\sin\theta)$. A single
$A$-edit with $\Delta a=a_s-a_r$ moves $u_a$ by $\Delta a\cos^2\theta$ and $e_A$ by $\Delta a\sin\theta\cos\theta$; the
readout recombines the two into $a_s$ exactly, and $|e_A|\le 2\sin\theta\cos\theta<\tau$ keeps the last unit silent.
Every single query therefore passes, although every edit that changes a variable misses its target. A joint edit
with $|\Delta a|+|\Delta b|\ge 3$ switches the last unit on and is wrong by at least 0.5. For any
$0^\circ<\theta<90^\circ$ (we use $\theta=\ExTheta^\circ$), none of the \ExSingleQueries{} single queries fails,
while \ExJointFailures{} of the \ExJointQueries{} joint queries do (\ExJointShare\%), all of them in the joint panel
(\ExPanelShare\% of it). If only $v_A$ is tilted and $v_B=(0,1,0,0)$, the $B$-edits land on their targets and no
joint query fails, as \cref{prop:redundant} requires.

\paragraph{How far a missed target can move the output.}
\Cref{prop:redundant} needs the edits of one variable to land exactly. A quantitative version says how far they may
miss. For a state $w$ at the site and a query that predicts class $c$, let the \emph{margin} $\mu_c(w)$ be the logit
of $c$ minus the largest other logit, so that the query passes when $\mu_c(w)>10^{-4}$. Write $r'=(a_r,b_t,n_r)$ and
$r''=(a_s,b_r,n_r)$ for the recipients of the two single queries that a joint query $(r,s,t)$ reduces to. By
\cref{lem:composition},
\begin{equation}\label{eq:split}
  J(r,s,t)=\SA(z_{r'},s)+u_A=\SB(z_{r''},t)+u_B,
\end{equation}
where $u_A=\bigl(\SB(z_r,t)-z_{r'}\bigr)(I-P_A)$ and $u_B=\bigl(\SA(z_r,s)-z_{r''}\bigr)(I-P_B)$: the joint state is
the state of either reduced single query, moved by the part of the other edit's miss that lies outside the reduced
query's own subspace.

\begin{proposition}\label{prop:margin}
The joint query $(r,s,t)$ and the two single queries $(r',s)$ and $(r'',t)$ predict the same class $c$. Let
$\mu_A=\mu_c(\SA(z_{r'},s))$ and $\mu_B=\mu_c(\SB(z_{r''},t))$ be the margins of the two single queries and
$\Delta_A=\mu_c(J(r,s,t))-\mu_A$ the change from the first to the joint query.
\begin{enumerate}[label=(\alph*)]
  \item If $\mu_c$ is differentiable at $\SA(z_{r'},s)$, then $\Delta_A=\nabla\mu_c(\SA(z_{r'},s))\cdot u_A+o(\|u_A\|)$.
  \item If $\mu_c$ is Lipschitz with constant $L$, then $|\Delta_A|\le L\|u_A\|$, and the joint query passes whenever
  $\mu_A-L\|u_A\|>10^{-4}$.
\end{enumerate}
The same holds for $\Delta_B=\mu_c(J(r,s,t))-\mu_B$ and $u_B$. Hence a joint query whose two single queries pass can
fail only if $\Delta_A<-\mu_A-10^{-4}$ and $\Delta_B<-\mu_B-10^{-4}$: in both orderings, the other edit's miss must
lower the margin by more than the single query's own margin.
\end{proposition}

\begin{proof}
All three queries predict the output of $(a_s,b_t,n_r)$. By \cref{eq:split}, $\mu_c(J(r,s,t))=\mu_c(\SA(z_{r'},s)+u_A)$,
so (a) is Taylor's theorem along $u_A$ and (b) the Lipschitz property. The joint query fails if
$\mu_c(J(r,s,t))<-10^{-4}$, which is $\Delta_A<-\mu_A-10^{-4}$, and equally $\Delta_B<-\mu_B-10^{-4}$.
\end{proof}

\Cref{prop:redundant} is the case $u_A=0$. \Cref{prop:margin} turns the check into a number that can be computed for
every joint query from the two single queries it reduces to. Only a drop of margin can make a joint query fail where
its single queries pass, so the score is the first-order estimate of the smaller of the two relative drops,
\[
  \sigma(r,s,t)=\min\Bigl(-\frac{\nabla\mu_c\cdot u_A}{\mu_A},\ -\frac{\nabla\mu_c\cdot u_B}{\mu_B}\Bigr),
\]
with each gradient taken at the state of its single query. An irreducible failure needs both actual relative drops to
exceed 1, and the rule $\sigma>1$ flags the joint queries for which both first-order estimates do. The unsigned score
$\rho$, the smaller of $|\nabla\mu_c\cdot u_A|/\mu_A$ and $|\nabla\mu_c\cdot u_B|/\mu_B$, measures the size of the
change whatever its direction. For an MLP readout $g(w)=\mathrm{relu}(wW_2+b_2)W_3+b_3$, the difference between the
logits of classes $c$ and $c'$ is Lipschitz with constant $\|W_2\|_2\,\|W_3(e_c-e_{c'})\|$, so (b) holds with
$L=\|W_2\|_2\max_{c\ne c'}\|W_3(e_c-e_{c'})\|$. LayerNorm leaves transformers without a useful global constant, and
for them only (a) applies.

\paragraph{What the score costs.}
The score needs the margin and its gradient at each single query, and both serve every joint query built on that single query: in the domains of Studies~1--3, each of the 1,458 single queries is one of the two single queries of 27 joint queries. Given the
single queries' margins, which a single test computes anyway, scoring all \ExJointQueries{} joint queries therefore
takes 1,458 gradients, which we charge at one model call each; the vectors $u_A$ and $u_B$, which come from natural
states and projectors without running the network; and two dot products of length $d$ per joint query,
\CostDotProducts{} in all. Running the joint queries instead takes \CostJointCalls{} model calls at four per query. The
8,748 queries of the joint panel need only the 972 single queries of the full-domain panel and their gradients.

\begin{remark}[Variables at different sites]\label{rem:sites}
Suppose $A$ is aligned at a site upstream of $B$'s, with natural states $z^A_x$ and $z^B_x$ at the two sites. A single
$A$-query edits $A$'s site and lets the network run on; a joint query $(r,s,t)$ does the same and then replaces the
$B$-component of the state $w$ that reaches $B$'s site, giving $w+(z^B_t-w)P_B$. If the $A$-edit lands on its
target, $w=z^B_{r''}$ and the joint query is the single $B$-query $(r'',t)$, as in \cref{prop:redundant}. There is no
reduction through $B$, because the $B$-edit acts on a state that the $A$-edit has already changed, and only one form
of \cref{eq:split} survives: the joint state is $\SB(z^B_{r''},t)+(w-z^B_{r''})(I-P_B)$. A joint failure whose single
queries $(r'',t)$ and $(r,s)$ both pass therefore needs only the $A$-edit to miss its target. Study~3
(\cref{sec:multisite}) asks whether this makes joint-only failures more common.
\end{remark}

\paragraph{What this means for testing.}
\Cref{prop:redundant} gives a condition under which joint tests are redundant that can be checked directly: do the
single edits land on their targets, up to directions the downstream computation ignores? For exact alignments, as
in Study~1, the check settles the question. Approximate alignments essentially never reproduce natural states
exactly. For them \cref{prop:margin} replaces the check with a margin condition: a joint query can fail only where
the other edit's miss is large compared with the margins of both single queries it reduces to. How often that happens,
and whether $\sigma$ finds the queries where it does, are empirical questions, which Studies~2--4 address.
